% SPDX-License-Identifier: Apache-2.0
\section{Words to a Frame and Back}
\label{sec:x-ethdma}

\S\ref{sec:x-dma} and \S\ref{sec:x-dmaw} move 32-bit words, because
that is what the bus moves. \S\ref{sec:x-ethslots} counts bytes,
because that is what a frame is. The two units here are that
difference and nothing else: neither touches the bus, and neither
knows where in memory the frame is.

\begin{center}\footnotesize
\begin{tabular}{@{}l@{}}
\code{words -- FrameOut --bytes-- EthTx --wire-- EthRx --bytes--
FrameIn -- words} \\
\end{tabular}
\end{center}

The run sends two frames out through the real transmitter, loops the
wire back, and checks what the receiver's bytes pack into. Putting
the MAC in the middle rather than joining the two adapters directly
is what makes the check worth running: the bytes that come back have
been framed, padded, checked and stripped, so the two units are
tested against the thing they will actually be joined to.

\subsection{Why the length is carried rather than counted}

A word says nothing about how many of its bytes are real. On the way
out the count is \code{tx\_length}, which the driver wrote. On the
way in it cannot be counted as the bytes go past, because the engine
that writes them to memory is told how many bytes to write when it
starts, and by then the bytes have not been seen.

So the receiver says. It takes a whole frame and checks it before
offering any of it, so the length is settled before the first byte
is offered, and this change gives it a port that says so. It reads
zero at every other time rather than holding the last frame's
length, so that a reader cannot mistake a stale length for a
current one.

\subsection{Two frames, not one}

The first version of this run sent one frame and proved less than it
looked.

A \code{FrameOut} that read whole words rather than stopping at the
byte count still marks the frame's last byte in the right place. The
frame on the wire is the right length, the words come back right, the
receiver reports the right count, and every check passes. The three
bytes it read past the frame do not lengthen that frame: they stay in
the transmitter, which has seen no last byte for them, and they go
out in front of the \emph{next} frame.

That fault was put into the unit deliberately to see whether the run
would catch it. With one frame it did not. With two it fails on the
second frame's first word and every word after it, because the whole
frame has moved along by three bytes.

\begin{center}\footnotesize
\begin{tabular}{@{}lll@{}}
\toprule
& one frame & two frames \\
\midrule
\code{FrameOut} reads past the count & passes & fails \\
\code{FrameIn} misaligns a partial word & fails & fails \\
\bottomrule
\end{tabular}
\end{center}

The lengths are 101 and 67 bytes. Both are above the 60 the
transmitter pads to, so no padding is in the way, and neither
divides by four: 101 leaves one real byte in its last word and 67
leaves three, so two of the three partial widths are covered rather
than the same one twice.

\lstinputlisting[rangebeginprefix=//\ begin\{,rangebeginsuffix=\},%
rangeendprefix=//\ end\{,rangeendsuffix=\},includerangemarker=false,%
linerange=in-in,caption={\code{lib/parts/src/ethdma.rs}: the
receiving half, which packs bytes into words and says how long the
frame was.},label={lst:x-ethdma}]{lib/parts/src/ethdma.rs}

\subsection{Two faults in the tools, found by simulating the netlist}

The Rust run passed before either netlist did, which is what the
netlist simulation is for.

A one-bit register incremented lowers to \code{slot + '1'}, and
\code{'1'} is a character rather than an unsigned, so \code{nvc}
refuses it. At every other width the same expression lowers to
\code{to\_unsigned(1, N)} correctly. Nothing in the tree had
incremented a one-bit register before, because a one-bit register is
usually a flag and a flag is set rather than counted. That is issue
461, and the slot here is written as a toggle instead.

A testbench reads the trace by port name, and both units had a port
called \code{out}. The second bound to the first's recorded signal
and failed in the generated VHDL, complaining about the length of a
string literal, 41 elements against 35: six hex digits, which is 24
bits, which is 32 less 9. Two more ports collided in the same run,
\code{bytes} and \code{go}, and neither was reported, because a
collision is only noticed when the widths differ. That is issue 462,
and it is the worse half: a silent collision is a co-simulation that
passes for the wrong reason. The ports here are named apart.

\lstinputlisting[style=ascii,caption={What \code{ex\_ethdma} printed
when the build ran it.},%
label={lst:x-ethdma-out}]{out_ethdma.txt}
